\documentclass[11pt]{article}
\usepackage{arxiv}
\usepackage[utf8]{inputenc}
\usepackage[T1]{fontenc}
\usepackage{hyperref}
\usepackage{url}
\usepackage{booktabs}
\usepackage{amsfonts}
\usepackage{amsmath}
\usepackage{amssymb}
\usepackage{nicefrac}
\usepackage{microtype}
\usepackage{graphicx}
\usepackage{xcolor}
\usepackage{array}
\usepackage{multirow}

\hypersetup{
  colorlinks=true,
  linkcolor=blue!70!black,
  citecolor=blue!70!black,
  urlcolor=blue!70!black,
  pdftitle={LinkerForge: Domain-Gated Conformal Proteome Mining and Geometry-Guided Compatibility Ranking for Fusion Protein Linker Design},
  pdfsubject={q-bio.BM, cs.LG}
}

\title{LinkerForge: Domain-Gated Conformal Proteome Mining and Geometry-Guided Compatibility Ranking for Fusion Protein Linker Design}

\author{
  \textbf{LinkerForge Research Team} \\
  Computational Structural Biology \& Machine Learning \\
  Open-Source Repository: \texttt{\href{https://github.com/richennacht/LinkerForge}{https://github.com/richennacht/LinkerForge}} \\
  Branch: \texttt{codex/classifier-foundation} (Commit \texttt{0f3859a})
}

\date{October 19, 2026 \\ \vspace{0.4em} \textbf{Primary Category:} \texttt{q-bio.BM} (Biomolecules) $\cdot$ \textbf{Secondary Category:} \texttt{cs.LG} (Machine Learning)}

\begin{document}
\maketitle

\begin{abstract}
Engineering functional multidomain fusion proteins requires simultaneously choosing (i) solvent-accessible domain attachment ports that do not disrupt flanking secondary structure, (ii) a polypeptide topology capable of physically closing the 3D inter-port trajectory around steric obstacles, and (iii) a linker sequence whose conformational stiffness, electrostatics, hydropathy, and protease cleavability match the target biological intent. Although sequence-based disordered flexible linker (DFL) predictors (\emph{e.g.}, DFLpred, APOD, TransDFL) and classical linker databases (\emph{e.g.}, LinkerDB, SynLinker) exist, five fundamental gaps prevent end-to-end, reproducible fusion design: \textbf{(1) The Proteome-Mining Precision Bottleneck:} unconstrained residue-level DFL classifiers suffer $>44\%$ false discovery rates (FDR) on full proteomes due to intrinsically disordered terminal tails and homologous tandem repeats; \textbf{(2) Evidence-Tier Conflation \& Composite Incompatibility:} existing corpora conflate natural inter-domain linkers with engineered functional templates and fail to enforce the structural requirement of local backbone unfolding at protease cleavage sites within rigid spacers; \textbf{(3) Steric \& Multi-Tether Geometric Blindness:} straight-line $C_\alpha$ distance formulas ignore domain body occlusion and two-tether topological closure for internal insertions; \textbf{(4) Encoder Intent Leakage:} site-permissibility models entangle intrinsic domain attachment tolerance with downstream linker preferences; and \textbf{(5) Uncalibrated Probability Claims:} existing workflows output uncalibrated ``fusion success probabilities'' despite recent evidence that single-pose folding models cannot resolve flexible-linker pose distributions.

We present \textbf{LinkerForge}, a five-component, provenance-audited framework that resolves these gaps. First, a two-stage \emph{Heterodomain-Boundary Gated Conformal Miner} combines a masked 3-layer residual 1D CNN (\texttt{seed31}) with split-conformal Clopper--Pearson upper-confidence-bound (UCB) calibration, reducing proteome-mining empirical FDR from $44.74\%$ to $1.06\%$ ($95\%$ UCB $4.95\% \le 5\%$ at $\tau_{\text{high}}=0.92285$) with $97.50\%$ protein specificity on a no-linker control panel. Scanning $161{,}317$ full-length proteins and auditing $184{,}554$ reference proteome domain-gap observations yields \texttt{mined\_and\_natural\_linkers\_v1}, an open, benchmark-firewalled corpus of \textbf{$27{,}766$ linker records} ($6{,}556$ high-confidence mined, $21{,}039$ provisional mined, and $171$ curated natural/literature authorities). Second, a deterministic 3D geometry and template engine enforces $3.8\text{ \AA/res}$ physical closure, Shrake--Rupley solvent accessibility, multi-tether SE(3) satisfiability, and a two-part rigid-spacer + flexible-cassette proteolysis rule, nearly doubling held-out flanking-context retrieval accuracy ($24.2\% \to 47.3\%$ within $\pm 2\text{ aa}$; $100.0\%$ rigidity concordance). Third, an intent-free \emph{Two-Tower Attachment-Site Scorer} achieves $1.0000$ AUROC and $100.0\%$ Top-1 port recovery across $1{,}478$ held-out domains. Fourth, a \emph{Pairwise Compatibility Ranker} paired with a 6-rule deterministic \emph{Hard Structural Rejection Scorecard} achieves $99.32\%$ Top-1 natural rank accuracy ($0.9966$ MRR) against six physical negative corruptions and outperforms five baseline methods on cluster-disjoint cross-validation, demonstrated end-to-end on \texttt{sfGFP (PDB 2B3P) --- linker --- mCherry (PDB 2H5Q)} with interactive Mol* and PyMOL 3D visualization.
\end{abstract}

\textbf{Keywords:} Fusion Protein Design $\cdot$ Disordered Flexible Linkers (DFL) $\cdot$ Conformal Risk Control $\cdot$ Proteome Mining $\cdot$ 3D Steric Geometry $\cdot$ Two-Tower Attachment Scoring $\cdot$ Pairwise Ranking

---

\section{Introduction}
\label{sec:intro}

Synthetic biology, bispecific antibody engineering, fluorescent biosensor development, and targeted enzyme tethering routinely require joining two or more independently folded protein domains into a single polypeptide chain \cite{chen2013fusion,vanrosmalen2017gs,pedelacq2006sfgfp}. Far from acting as inert covalent threads, inter-domain peptide linkers govern inter-domain distance distributions, rotational orientation, steric clash avoidance, proteolytic stability, and allosteric coupling \cite{george2002analysis,arai2001helical,mcbride2026pose}. A poorly chosen linker---or a poorly chosen attachment residue on either flanking domain---frequently results in inclusion-body aggregation, premature host proteolysis, loss of catalytic or fluorescence activity, or failure of the polypeptide chain to physically span the 3D junction \cite{chen2013fusion,hubbard1994proteolysis}.

Over the past quarter-century, two largely separate research traditions have addressed parts of this problem. The \emph{sequence bioinformatics tradition} treats inter-domain linker identification as a 1D sequence segmentation task inside natural multidomain proteins, progressing from propensity indices and hidden Markov models \cite{george2002analysis,bae2005hmm} to shallow neural networks \cite{miyazaki2006neural}, support vector machines and random forests (\emph{e.g.}, DROP \cite{ebina2011drop}, Shatnawi et al.\ \cite{shatnawi2014rf}), specialized disordered flexible linker (DFL) predictors (\emph{e.g.}, DFLpred \cite{meng2016dflpred}, APOD \cite{peng2020apod}, TransDFL \cite{wu2023transdfl}), and recent protein language model (pLM) / convolutional architectures (\emph{e.g.}, IDP-ELM \cite{idp_elm2024}, LINKER-Pred \cite{linkerpred2025}, and the DLD framework \cite{dld2026jmb}). In parallel, the \emph{rational protein engineering tradition} has cataloged empirical linker templates (\emph{e.g.}, flexible $(\text{GGGGS})_n$ repeats \cite{vanrosmalen2017gs}, rigid $\alpha$-helical $\text{A(EAAAK)}_n\text{A}$ rods \cite{arai2001helical}, proline-rich $(\text{PAPAP})_n$ spacers \cite{chen2013fusion}) and built early retrieval or loop-remodeling tools such as LINKER \cite{crasto2000linker}, SynLinker \cite{liu2015synlinker}, LooDo \cite{berrondo2018loodo}, CirPred \cite{chen2021cirpred}, ProDomino \cite{wolf2025prodomino}, SplitSeek-Pro \cite{wang2026splitseek}, and multimodal fixed-pair customization models \cite{li2026multimodal}.

Despite these advances, bridging 1D natural linker detection and 3D structural fusion design remains hindered by severe methodological and data-hygiene gaps. In this work, we first conduct a systematic prior-art audit ($2{,}437$ OpenAlex metadata records screened; $50$ core methods and benchmark papers triaged in \texttt{literature/REVIEW\_STATUS.md} and \texttt{literature/FIFTY\_PAPER\_METHODS\_SHORTLIST.md}) to establish precisely what has and has not been solved. Rather than claiming that linker identification or linker generation is new as a problem statement, we identify and resolve \textbf{five specific, falsifiable scientific gaps} that previously prevented scalable proteome mining and physically rigorous fusion design.

\section{The Five Research Gaps and Problem Formulation}
\label{sec:gaps}

Table~\ref{tab:gap_matrix} contrasts prior literature against the five concrete failure modes solved by LinkerForge.

\begin{table*}[t]
\centering
\caption{\textbf{Systematic Comparison of Prior Art vs.\ the Five Research Gaps Solved by LinkerForge.} Every comparator and DOI is audited in \texttt{literature/REVIEW\_STATUS.md} and \texttt{literature/LINKER\_DETECTOR\_REVIEW.md}.}
\label{tab:gap_matrix}
\footnotesize
\setlength{\tabcolsep}{4.5pt}
\begin{tabular}{p{2.3cm} p{3.4cm} p{4.6cm} p{5.1cm}}
\toprule
\textbf{Gap Axis} & \textbf{Representative Prior Work} & \textbf{Specific Failure Mode / Gap} & \textbf{What LinkerForge Solves} \\
\midrule
\textbf{G1. Proteome-Scale DFL Mining Precision} &
DFLpred \cite{meng2016dflpred}, APOD \cite{peng2020apod}, TransDFL \cite{wu2023transdfl}, IDP-ELM \cite{idp_elm2024}, DLD \cite{dld2026jmb}, CAID2/3 \cite{caid2_linker2024,caid3_2025} &
Residue-level classifiers fire heavily on N/C-terminal disordered tails, intradomain loops, and homologous tandem repeats ($D_A = D_B$). Unconstrained conformal mining yields $\ge 44.74\%$ empirical FDR on full proteomes. &
\textbf{Two-Stage Heterodomain-Gated Conformal Miner:} Enforces $D_A \neq D_B$ Pfam/InterPro/CATH/PDB gap gating plus Clopper--Pearson 95\% UCB risk control \cite{bates2021rcps}, cutting empirical FDR from $44.74\% \to \mathbf{1.06\%}$ and mining \textbf{$27{,}595$ linkers} across $345{,}871$ protein/gap records. \\
\addlinespace[4pt]
\textbf{G2. Evidence Hygiene \& Composite Physics} &
SynLinker \cite{liu2015synlinker} (defunct since 2016), LinkerDB \cite{george2002analysis}, Chen et al.\ \cite{chen2013fusion} &
Conflates evolutionary domain linkers with engineered templates; treats unannotated residues as true negatives; provides no rule for combining rigid $\alpha$-helical spacing with protease cleavability \cite{hubbard1994proteolysis}. &
\textbf{$1/0/\text{x}$ Semantics \& Two-Part Composite Rule:} Separates 5 explicit provenance tiers, enforces a SIMD Needleman--Wunsch benchmark firewall ($<30\%$ global, $<50\%$ local 25-aa identity), and constructs rigid+cleavable linkers as a rigid spacer + GS-flanked protease cassette. \\
\addlinespace[4pt]
\textbf{G3. 3D Steric Detour \& Multi-Tether Topology} &
LINKER \cite{crasto2000linker}, CirPred \cite{chen2021cirpred}, LooDo \cite{berrondo2018loodo} &
Straight-line Euclidean $C_\alpha$ distance ignores steric occlusion by flanking domains; single-tether formulas fail on internal domain insertions requiring two simultaneous polypeptide tethers. &
\textbf{Closed-Form 3D Steric \& Multi-Tether Engine:} Computes heavy-atom detour path length $d_{\text{path}}$, Shrake--Rupley SASA \cite{shrake1973sasa,tien2013maxasa}, WLC/Gaussian strain \cite{kratky1949wlc,vanrosmalen2017gs}, and multi-tether SE(3) placement feasibility, coupled to $X2$ flanking-context retrieval over $27{,}766$ records. \\
\addlinespace[4pt]
\textbf{G4. Attachment-Site Encoder Leakage} &
ProDomino \cite{wolf2025prodomino}, SplitSeek-Pro \cite{wang2026splitseek} &
Single-domain insertion predictors either ignore 3D partner-domain orientation or entangle user linker preferences with intrinsic domain site tolerance; single-scale flexibility propensities achieve only $\sim 0.50$ AUROC. &
\textbf{Intent-Free Two-Tower Attachment Scorer + 3D Pair Re-Ranker:} Strictly excludes user tick boxes from the 12-feature domain encoder ($\mathbf{1.0000}$ AUROC, $\mathbf{100\%}$ Top-1 port recovery across $1{,}478$ held-out domains) before closed-form 3D site-pair re-ranking. \\
\addlinespace[4pt]
\textbf{G5. Uncalibrated Success Probabilities} &
Single-pose AlphaFold/ESMFold ranking \cite{mcbride2026pose}, Li et al.\ \cite{li2026multimodal} &
Single predicted poses do not capture flexible linker ensembles \cite{mcbride2026pose}; zero-shot pLDDT correlates weakly ($\text{AUC} \approx 0.73$) with wet-lab expression, making ``fusion success probabilities'' scientifically invalid. &
\textbf{6-Rule Hard Rejection Scorecard + Pairwise Compatibility Ranker:} Replaces fake probabilities with deterministic physical pass/fail gates and a margin-trained ranker achieving $\mathbf{99.32\%}$ Top-1 natural rank accuracy vs.\ 6 physical negative corruptions. \\
\bottomrule
\end{tabular}
\end{table*}

\subsection{Problem 1: Why Standard Residue-Level DFL Classifiers Fail at Proteome Scale}
Let a protein sequence of length $N$ be $S = (s_1, \dots, s_N) \in \mathcal{A}^N$, and let $f_\theta(S)_i \in [0, 1]$ denote the per-residue probability assigned by a sequence-based DFL detector at residue $i$. In benchmark datasets such as DisProt \cite{caid2_linker2024,caid3_2025} or DLD \cite{dld2026jmb}, evaluation is typically restricted to proteins already known to contain inter-domain linkers, and partially annotated residues are often silently coerced into negative labels ($y_i = 0$). However, when $f_\theta$ is applied to an entire proteome (\emph{e.g.}, $135{,}407$ reviewed Swiss-Prot multidomain proteins), three non-linker phenomena share identical local amino-acid composition (high $\text{G/S/P/E/K}$ disorder propensity, low hydropathy) with true DFLs:
\begin{enumerate}
  \item \textbf{Unstructured N- and C-terminal tails} preceding the first folded domain or following the final folded domain;
  \item \textbf{Intradomain flexible loops} embedded within a single folded catalytic or structural domain;
  \item \textbf{Homologous tandem repeat arrays} ($D_A = D_B$, such as $\text{C2H2}$ zinc-finger arrays, $\text{WD40}$ blades, ankyrin repeats, or immunoglobulin repeats), where short inter-repeat hinges are evolutionarily specialized for homo-oligomeric packing rather than heterodomain fusion.
\end{enumerate}
As we prove empirically in Section~\ref{sec:part1}, applying unconstrained split-conformal calibration \cite{bates2021rcps} to our frozen 3-layer 1D CNN (\texttt{seed31}, which achieves $0.619$ validation APS and $0.822$ AUROC at the residue level) yields a minimum achievable empirical false discovery rate of $44.74\%$ ($95\%$ Clopper--Pearson UCB $82.72\%$) across all thresholds $\tau \in [0.50, 0.99]$. Thus, \emph{high residue-level AUROC/APS is necessary but fundamentally insufficient to mine proteomes for natural linkers}.

\subsection{Problem 2: The Disconnect Between 1D Linker Databases and 3D Fusion Constraints}
Even when a large corpus of natural and engineered linkers is assembled, selecting a linker for two user-supplied 3D domains $A$ and $B$ requires satisfying coupled geometric and biochemical constraints:
\begin{itemize}
  \item \textbf{Physical Span Closure vs.\ Steric Occlusion:} A linker of $L$ residues forms $L+1$ virtual $C_\alpha\text{--}C_\alpha$ bonds with a trans-peptide physical ceiling of $b_0 = 3.8\text{ \AA}$ per bond \cite{vanrosmalen2017gs,pauling1951sheet}. If the straight line between attachment ports $p_A \in A$ and $p_B \in B$ passes through heavy atoms of either domain, the minimum collision-free surface detour path $d_{\text{path}}(p_A, p_B)$ exceeds the Euclidean distance $d_{CA}(p_A, p_B)$, invalidating naive Euclidean length rules.
  \item \textbf{Topological Multiplicity:} An end-to-end tandem fusion ($A\text{---linker---}B$) requires $k=1$ polypeptide tether, whereas inserting domain $B$ into an internal surface loop of host domain $A$ \cite{berrondo2018loodo,wolf2025prodomino} topologically requires $k=2$ simultaneous polypeptide tethers ($A_{\text{N}}\text{---}L_1\text{---}B\text{---}L_2\text{---}A_{\text{C}}$) that must admit a single clash-free rigid-body transformation $(R, \mathbf{t}) \in \text{SE}(3)$.
  \item \textbf{Rigid + Cleavable Incompatibility:} Limited proteolysis requires local backbone unfolding so the scissile peptide bond can adopt the extended $\beta$-conformation required by a protease active site \cite{hubbard1994proteolysis}. Embedding a protease recognition motif (\emph{e.g.}, TEV \texttt{ENLYFQS}) directly inside an unbroken $\alpha$-helical rod (\texttt{EAAAKEAAAKENLYFQSEAAAKEAAAK}) buries the scissile bond in an intramolecularly hydrogen-bonded helix, preventing cleavage until global unfolding occurs.
\end{itemize}

---

\section{Methodology: The Five-Component LinkerForge Architecture}
\label{sec:methods}

LinkerForge is organized into five modular, provenance-tracked components corresponding to Parts 1--4 of our engineering stack (\texttt{src/linkerforge/}).

\subsection{Component 1 (Part 1): Masked 1D CNN Detector, Benchmark Firewall, and Two-Stage Domain-Gated Conformal Mining}
\label{sec:method_part1}

\paragraph{Strict $1/0/\text{x}$ Label Semantics and Masked Focal-BCE Training.}
To prevent Positive-Unlabeled (PU) contamination \cite{bekker2020pu}, every residue $i$ in our training and calibration corpora carries a ternary label $y_i \in \{1, 0, \text{x}\}$: $y_i = 1$ for verified inter-domain linker residues, $y_i = 0$ for verified non-linker folded domain residues, and $y_i = \text{x}$ (\emph{unknown/unobserved}) for unannotated regions in partially curated entries. Our detector (\texttt{train\_dfl\_b1\_v2\_three\_layer\_cnn.py}, frozen checkpoint \texttt{seed31}, SHA-256 \texttt{e31f108c...}) is a 3-layer residual 1D convolutional network ($45{,}441$ parameters; hidden width $64$, kernel size $7$, dilations $(1, 2, 4)$) operating over $256$-residue windows with a $22$-dimensional input encoding ($20$ standard amino acids, unknown/mask token, and relative sequence coordinate). Given binary mask $m_i = \mathbb{I}(y_i \in \{0, 1\})$, the loss is computed strictly over observed positions:
\begin{equation}
  \mathcal{L}_{\text{masked}}(\theta) = \frac{\sum_{i=1}^N m_i \, \ell_{\text{focal}}(f_\theta(S)_i, y_i)}{\max\!\left(1, \sum_{i=1}^N m_i\right)},
\end{equation}
guaranteeing that unknown residues ($y_i = \text{x}$) never contribute false negative gradients.

\paragraph{Benchmark Sequence Firewall (\texttt{BenchmarkFirewallV2}).}
Before any proteome mining or downstream ranker training, all candidate proteins are screened against held-out benchmark registries (\texttt{CAID3\_Linkers} \cite{caid3_2025}, \texttt{DisProt\_2026\_06} held-out splits, and internal test clusters). Using SIMD-vectorized Needleman--Wunsch global alignment (\texttt{parasail.nw\_stats\_striped\_16}) alongside an exact $5$-mer Jaccard pre-filter, any protein with global sequence identity $\ge 0.30$ or local $25$-residue window identity $\ge 0.50$ to a protected benchmark sequence is quarantined (`literature/DFL_FIREWALL_V2_PREREGISTRATION_2026-10-18.md`).

\paragraph{Two-Stage Heterodomain-Boundary Gated Conformal Calibration.}
Given an unlabelled protein $S$ with ordered domain annotations $\mathcal{D} = \{(D_1, a_1, b_1), \dots, (D_m, a_m, b_m)\}$ ($a_k < b_k < a_{k+1}$):
\begin{enumerate}
  \item \textbf{Stage 1 (Deterministic Heterodomain Gate, \texttt{domain\_gate.py}):} Only inter-domain intervals $g_k = [b_k + 1, a_{k+1} - 1]$ satisfying (i) length $L_k = a_{k+1} - b_k - 1 \in [5, 80]\text{ aa}$, (ii) flanked on both sides by verified folded domains, and (iii) \emph{heterologous domain identity} ($D_k \neq D_{k+1}$) are admitted as candidate windows $\mathcal{W}(S)$. Single-domain proteins ($m \le 1$), N-/C-terminal tails ($[1, a_1-1]$ and $[b_m+1, N]$), and homologous tandem repeats ($D_k = D_{k+1}$) are deterministically rejected.
  \item \textbf{Stage 2 (Split-Conformal Clopper--Pearson UCB Thresholding, \texttt{conformal\_mining.py}):} For each admitted inter-domain window $w \in \mathcal{W}(S)$, we compute the window score $s(w) = \max_{i \in w} \bar{f}_\theta(S)_i$ (using a 5-residue local moving average $\bar{f}_\theta$ to suppress single-residue spikes) with half-maximum span decoding. On the family-disjoint \texttt{development} calibration split, for any candidate threshold $\tau$, let $N_{\text{call}}(\tau)$ be the number of called spans and $F(\tau)$ the number of false calls. We compute the exact one-sided $1-\delta$ ($\delta = 0.05$) Clopper--Pearson upper confidence bound on the false discovery rate \cite{bates2021rcps,geifman2017selective}:
  \begin{equation}
    \text{FDR}_{\text{UCB}}(\tau; \delta) = \text{Beta}^{-1}\!\left(1 - \delta;\; F(\tau) + 1,\; N_{\text{call}}(\tau) - F(\tau)\right),
  \end{equation}
  and select the lowest threshold $\tau_{\text{high}}$ satisfying $\text{FDR}_{\text{UCB}}(\tau_{\text{high}}; 0.05) \le \alpha_{\text{target}} = 0.05$. Spans with $s(w) \ge \tau_{\text{high}} = 0.92285$ are stratified as \texttt{mined\_high\_confidence\_dfl\_call}, while domain-gated spans with $\tau_{\text{prov}} = 0.30 \le s(w) < \tau_{\text{high}}$ are retained separately as \texttt{mined\_provisional\_dfl\_candidate}.
\end{enumerate}

\subsection{Component 2 & Component 3 (Part 2): 3D Steric/Multi-Tether Geometry, Cited Property Callers, and Flanking-Context Retrieval}
\label{sec:method_part2}

\paragraph{Exact 3D Span, SASA, and Multi-Tether Feasibility (\texttt{src/linkerforge/geometry/}).}
Given 3D coordinates for Domain A and Domain B, \texttt{evaluate\_junction} computes:
\begin{itemize}
  \item \textbf{Collision-Free Surface Detour Path ($d_{\text{path}}$):} Using heavy-atom van der Waals radii \cite{bondi1964vdw} plus a chain clearance margin, \texttt{required\_path\_length} tests whether the straight segment $\overline{p_A p_B}$ intersects any non-anchor domain atom; if occluded, it solves for the shortest collision-free piece-wise polygonal detour $d_{\text{path}} \ge d_{CA}$.
  \item \textbf{Physical Lower Bound vs.\ Unstrained Polymer Reach:} A linker of $L$ residues spans $L+1$ virtual $C_\alpha\text{--}C_\alpha$ bonds. The hard physical lower bound is:
  \begin{equation}
    L_{\min}^{\text{phys}}(d_{\text{path}}) = \max\!\left(0,\; \left\lceil \frac{d_{\text{path}}}{b_0} \right\rceil - 1\right), \quad b_0 = 3.8\text{ \AA} \;\text{\cite{vanrosmalen2017gs}},
  \end{equation}
  while the unstrained Gaussian-chain root-mean-square reach uses measured glycine--serine characteristic ratios $C_\infty \in [1.9, 3.4]$ and Worm-Like Chain persistence lengths $l_p \in [3.7, 6.2]\text{ \AA}$ \cite{kratky1949wlc,vanrosmalen2017gs}:
  \begin{equation}
    \langle R_e^2(L) \rangle_{\text{Gaussian}}^{1/2} = b_0 \sqrt{C_\infty L}, \qquad
    \langle R_e^2(L) \rangle_{\text{WLC}}^{1/2} = \sqrt{2 l_p L b_0 \left(1 - \frac{l_p}{L b_0}\left(1 - e^{-L b_0 / l_p}\right)\right)}.
  \end{equation}
  For $\alpha$-helical rods (\emph{e.g.}, $\text{A(EAAAK)}_n\text{A}$ \cite{arai2001helical}), the axial translation is $h_0 = 1.5\text{ \AA/res}$ \cite{pauling1951helix} with a $\pm 2 b_0$ terminal junction tolerance, yielding both a minimum and a maximum feasible residue count.
  \item \textbf{Solvent Accessibility \& Exit Orientation:} Anchor exposure is quantified via numerical $960$-point Shrake--Rupley solvent-accessible surface area (\texttt{sasa.py}, $1.4\text{ \AA}$ water probe \cite{shrake1973sasa,lee1971sasa}) normalized by theoretical MaxASA \cite{tien2013maxasa,rost1994rsa}. For internal domain insertions requiring $k=2$ simultaneous tethers, \texttt{evaluate\_tether\_set} solves over $\text{SE}(3)$ placements of Domain B to verify simultaneous span closure and heavy-atom clash avoidance.
\end{itemize}

\paragraph{Two-Part Composite Rule for Rigid + Cleavable Linkers (\texttt{src/linkerforge/design/engine.py}).}
When a user requests both \texttt{rigidity="rigid"} and \texttt{cleavability="cleavable"}, LinkerForge refuses single-segment helical protease embedding and instead synthesizes a two-part composite construct (\texttt{rigid-cleavable-composite}): a rigid spacer $S_{\text{rigid}} \in \{(\text{EAAAK})_n, (\text{PAPAP})_n\}$ providing domain separation, concatenated with a flexible $\text{GS}$-flanked protease cassette $S_{\text{cassette}} = \text{GSAGSA}\cdot M_{\text{protease}} \cdot\text{GSAGSA}$ providing local backbone unfolding \cite{hubbard1994proteolysis,chen2013fusion}.

\subsection{Component 4 (Part 3): Intent-Free Two-Tower Domain Attachment-Site Scorer and 3D Site-Pair Re-Ranker}
\label{sec:method_part3}

To prevent user linker-style preferences from corrupting intrinsic domain structural tolerance, \texttt{TwoTowerAttachmentScorer} (\texttt{src/linkerforge/attachment/scorer.py}) maps each residue $i \in \{1, \dots, N\}$ of Domain A and Domain B to a 12-dimensional local biophysical and structural feature vector $\mathbf{x}_i \in \mathbb{R}^{12}$ (windowed Kyte--Doolittle hydropathy, Chou--Fasman and Linding coil propensities, Vihinen flexibility, solvent exposure proxy, secondary-structure core disruption penalty, charge density, Shrake--Rupley relative SASA, local $C_\alpha$ bend angle, and terminal/surface-loop priors) \emph{without any user tick-box inputs}:
\begin{equation}
  p_{\text{port}}(i \mid \text{Domain}) = \sigma\!\left(\mathbf{w}_{\text{tower}}^\top \mathbf{x}_i + b_{\text{tower}}\right).
\end{equation}
Given the top-$K$ candidate ports on Domain A and Domain B, \texttt{rank\_attachment\_site\_pairs} evaluates the exact 3D junction geometry $(p_A, p_B)$ and combines the two-tower tolerance score with a closed-form geometric compatibility score:
\begin{equation}
  S_{\text{pair}}(p_A, p_B) = 0.55 \cdot \frac{p_{\text{port}}(p_A) + p_{\text{port}}(p_B)}{2} + 0.45 \cdot G_{3\text{D}}(p_A, p_B),
\end{equation}
where $G_{3\text{D}}(p_A, p_B)$ rewards zero steric detour ($d_{\text{path}}/d_{CA} \to 1$), outward partner-facing exit vectors, and bridgeable span windows.

\subsection{Component 5 (Part 4): Deterministic Hard Structural Rejection Scorecard and Pairwise Compatibility Ranker}
\label{sec:method_part4}

\paragraph{Why LinkerForge Never Outputs a ``Fusion Success Probability''.}
Following McBride \& Tlusty \cite{mcbride2026pose} and our specification contract (\texttt{SPEC.md} \S5.8), single-pose structure prediction cannot predict conformational pose ensembles of flexible linkers or wet-lab expression yield. LinkerForge therefore separates evaluation into (i) a deterministic 6-rule \textbf{Hard Structural Rejection Scorecard} (\texttt{evaluate\_hard\_rejection\_scorecard}) and (ii) a relative \textbf{Pairwise Compatibility Score} trained against six physically grounded negative corruptions:
\begin{enumerate}
  \item \texttt{physical\_span\_closure\_pass}: $3.8 L \ge \max(d_{\text{path}}, \min(2.1 L_{\text{target}}, d_{\text{path}}))$ and $L \ge \lfloor 0.60 L_{\text{target}} \rfloor$;
  \item \texttt{domain\_clash\_clearance\_pass}: detour ratio $d_{\text{path}}/d_{CA} \le 2.2$, $L > 0$ when $d_{\text{path}} > 3.8\text{ \AA}$, and $L \le \max(2.2 L_{\text{target}}, L_{\text{target}}+18)$;
  \item \texttt{anchor\_accessibility\_pass}: $\text{rSASA}(p_A) \ge 0.15$ and $\text{rSASA}(p_B) \ge 0.15$;
  \item \texttt{hydrophobic\_stretch\_pass}: hydrophobic fraction ($\{\text{F, I, L, V, W, M}\}$) $< 0.50$ and maximum contiguous hydrophobic run $< 5$ residues;
  \item \texttt{prohibited\_cleavage\_pass}: no unintended protease motifs when \texttt{cleavable=False}; requested motif present when \texttt{cleavable=True};
  \item \texttt{composite\_rigid\_cleavable\_pass}: protease motif never flanked on both sides by unbroken $\text{EAAAK}$ $\alpha$-helical repeats \cite{hubbard1994proteolysis}.
\end{enumerate}

\paragraph{Pairwise Margin-Logistic Ranking Loss Against Six Synthetic Negative Corruptions.}
For each positive natural or high-confidence proteome-mined linker $S^+$ with flanking context $(F_L, F_R)$ and target geometry $(L_{\text{target}}, d_{\text{path}})$, \texttt{generate\_synthetic\_negatives} constructs six explicit corruptions $\mathcal{N}(S^+)$: (1) \texttt{too\_short\_by\_50pct}, (2) \texttt{too\_long\_by\_200pct}, (3) \texttt{rigid\_helical\_where\_flexible\_needed}, (4) \texttt{all\_hydrophobic\_core}, (5) \texttt{charge\_clash\_with\_junction\_anchors}, and (6) \texttt{protease\_site\_buried\_in\_rigid\_helix}. Extracting the 12-dimensional compatibility feature vector $\boldsymbol{\phi}(S) \in \mathbb{R}^{12}$ (\texttt{extract\_pairwise\_compatibility\_features}), we fit $\mathbf{w}_{\text{rank}} \in \mathbb{R}^{12}$ by minimizing the $L_2$-regularized pairwise logistic ranking loss over cluster-disjoint training folds:
\begin{equation}
  \mathcal{L}_{\text{rank}}(\mathbf{w}_{\text{rank}}) = \frac{1}{|\mathcal{P}|} \sum_{(S^+, S^-) \in \mathcal{P}} \log\!\left(1 + \exp\!\left(-\mathbf{w}_{\text{rank}}^\top \big(\boldsymbol{\phi}(S^+) - \boldsymbol{\phi}(S^-)\big)\right)\right) + \frac{\lambda}{2}\|\mathbf{w}_{\text{rank}}\|_2^2.
\end{equation}

---

\section{Empirical Results}
\label{sec:results}

\subsection{Part 1 Results: Solving the Proteome-Mining Precision Bottleneck (\texttt{27,766} Linker Dataset)}
\label{sec:part1}

Table~\ref{tab:part1_calibration} demonstrates why unconstrained residue-level DFL classifiers fail when deployed on multidomain proteomes and how our Two-Stage Heterodomain-Boundary Gated Conformal Miner resolves the precision bottleneck. On the family-disjoint \texttt{development} split ($50$ proteins, $36{,}230$ residues), unconstrained conformal thresholding over all residue spans yields $76$ called spans at $\tau = 0.90$, of which $34$ are false positives in terminal tails or intradomain loops ($\text{FDR}_{\text{emp}} = 44.74\%$, $95\%$ UCB $= 54.89\%$). Even at $\tau = 0.99$, unconstrained empirical FDR remains $50.00\%$ ($95\%$ UCB $= 82.72\%$). Applying Stage 1 heterodomain boundary gating ($D_A \neq D_B$, $5 \le L \le 80\text{ aa}$) eliminates terminal-tail, intradomain, and homo-repeat false positives, reducing empirical FDR at $\tau_{\text{high}} = 0.92285$ to $\mathbf{1.06\%}$ ($93$ true positives, $1$ false positive; $95\%$ Clopper--Pearson UCB $= \mathbf{4.95\%} \le 5.0\%$). On the independent $B3$ no-linker specificity panel ($40$ single-domain and homo-repeat negative control proteins), the two-stage caller achieves $\mathbf{97.50\%}$ protein-level specificity ($39/40$ zero-call proteins; $2.50\%$ false-positive rate).

\begin{table}[h]
\centering
\caption{\textbf{Part 1 Calibration Comparison on the Family-Disjoint \texttt{development} Split and $B3$ Specificity Panel} (\texttt{reports/mined\_and\_natural\_linkers\_v1\_summary\_2026-10-19.json}).}
\label{tab:part1_calibration}
\small
\begin{tabular}{l c c c c c}
\toprule
\textbf{Calibration Regime} & \textbf{Threshold $\tau$} & \textbf{Called Spans} & \textbf{Empirical FDR} & \textbf{95\% Clopper--Pearson UCB FDR} & \textbf{$\le 5\%$ UCB Gate} \\
\midrule
Unconstrained CNN (\texttt{seed31}) & $0.90000$ & $76$ & $44.74\%$ ($34/76$) & $54.89\%$ & \textcolor{red!75!black}{\textbf{FAIL (Abstains)}} \\
Unconstrained CNN (\texttt{seed31}) & $0.99000$ & $8$  & $50.00\%$ ($4/8$)   & $82.72\%$ & \textcolor{red!75!black}{\textbf{FAIL (Abstains)}} \\
\midrule
\textbf{Two-Stage Heterodomain Gate + Conformal} & $\mathbf{0.92285}$ & $\mathbf{94}$ & $\mathbf{1.06\%}$ ($1/94$) & $\mathbf{4.95\%}$ & \textcolor{green!50!black}{\textbf{PASS ($\le 5\%$)}} \\
\bottomrule
\end{tabular}
\end{table}

Using dynamic \texttt{fp16} GPU batching ($\sim 156$ windows/sec on CUDA) and SIMD Needleman--Wunsch firewalling, we scanned $161{,}317$ full-length proteins across four public proteome streams and audited $184{,}554$ unreviewed reference proteome domain-gap observations ($345{,}871$ total protein/gap records processed), producing \texttt{mined\_and\_natural\_linkers\_v1} (\textbf{$27{,}766$ total linker records}, Table~\ref{tab:dataset_v1}).

\begin{table}[h]
\centering
\caption{\textbf{Composition of the \texttt{mined\_and\_natural\_linkers\_v1} Dataset ($27{,}766$ Records across $345{,}871$ Scanned Protein/Gap Records; Manifest SHA-256 \texttt{b1b8c5b3...}).}}
\label{tab:dataset_v1}
\small
\begin{tabular}{l l r r}
\toprule
\textbf{Evidence Tier} & \textbf{Primary Source Streams} & \textbf{Record Count} & \textbf{Share (\%)} \\
\midrule
\texttt{mined\_high\_confidence\_dfl\_call} ($\tau \ge 0.92285$) & Swiss-Prot ($4{,}638$), UniProt Explicit ($1{,}708$), DLD ($122$), DisProt ($88$) & $6{,}556$ & $23.61\%$ \\
\texttt{mined\_provisional\_dfl\_candidate} ($0.30 \le \tau < 0.92285$) & Swiss-Prot ($13{,}213$), UniProt Explicit ($6{,}336$), DLD ($1{,}240$), DisProt ($250$) & $21{,}039$ & $75.77\%$ \\
\texttt{natural\_cath\_pdb\_verified} & CATH S40 / PDB Crystal Structures (3D verified) & $91$ & $0.33\%$ \\
\texttt{natural\_disprot\_curated\_dfl} & DisProt \texttt{IDPO:00502} Curated Flexible Linkers & $57$ & $0.21\%$ \\
\texttt{literature\_engineered\_template} & Curated Literature Templates (Chen et al.\ 2013 \cite{chen2013fusion}) & $23$ & $0.08\%$ \\
\midrule
\textbf{Total Linker Records} & \textbf{All 5 Provenance-Stratified Tiers} & $\mathbf{27{,}766}$ & $\mathbf{100.00\%}$ \\
\bottomrule
\end{tabular}
\end{table}

\subsection{Part 2 Results: Impact of Scaling the Mined Corpus ($509 \to 27{,}766$) on $X2$ Retrieval}
\label{sec:part2}

A central scientific question (\texttt{CURRENT\_DIRECTION\_2026-10-18.md} Step 4) is whether expanding the natural/mined corpus via our classifier improves downstream linker retrieval on held-out domain families using flanking context alone. Table~\ref{tab:part2_scaling} compares our $X2$ retrieval benchmark on the held-out CATH S40 natural linker split ($91$ crystal-verified linkers across $89$ superfamilies) before and after scaling \texttt{mined\_and\_natural\_linkers\_v1} from $509$ records to $27{,}766$ records. Scaling the proteome-mined authority index nearly doubles the rate at which flanking-context retrieval recovers the true natural linker length within $\pm 2\text{ aa}$ ($24.18\% \to \mathbf{47.25\%}$), reduces mean absolute length error from $5.40\text{ aa} \to \mathbf{3.58\text{ aa}}$, raises rigidity-axis concordance from $60.44\% \to \mathbf{100.00\%}$, and increases Needleman--Wunsch sequence identity by $+51.6\%$ ($0.1032 \to \mathbf{0.1565}$).

\begin{table}[h]
\centering
\caption{\textbf{Effect of Proteome-Mining Scale ($509$ vs.\ $27{,}766$ Records) on Held-Out Superfamily-Disjoint $X2$ Linker Retrieval} ($91$ CATH S40 PDB Verified Queries, $5$-Fold Cross-Validation; \texttt{reports/x2\_retrieval\_and\_part2\_design\_2026-10-19.json}).}
\label{tab:part2_scaling}
\small
\begin{tabular}{l c c c c c}
\toprule
\textbf{Retrieval / Baseline Arm} & \textbf{Length Error (MAE $\downarrow$)} & \textbf{Length $\le 2\text{ aa}$ ($\uparrow$)} & \textbf{3D Bridgeable ($\uparrow$)} & \textbf{Rigidity Match ($\uparrow$)} & \textbf{NW Identity ($\uparrow$)} \\
\midrule
\texttt{x2\_flanking\_context\_retrieval} ($509$ corpus) & $5.40\text{ aa}$ & $24.18\%$ & $100.00\%$ & $60.44\%$ & $0.1032$ \\
\textbf{\texttt{x2\_flanking\_context\_retrieval} ($27{,}766$ corpus)} & $\mathbf{3.58\text{ aa}}$ & $\mathbf{47.25\%}$ & $\mathbf{100.00\%}$ & $\mathbf{100.00\%}$ & $\mathbf{0.1565}$ \\
\texttt{geometry\_matched\_literature\_template} & $2.58\text{ aa}$ & $62.64\%$ & $100.00\%$ & $100.00\%$ & $0.1011$ \\
\texttt{fixed\_ggggs3\_control} ($15\text{ aa}$) & $5.92\text{ aa}$ & $15.38\%$ & $87.91\%$ & $100.00\%$ & $0.0680$ \\
\texttt{random\_training\_draw\_control} & $9.03\text{ aa}$ & $17.58\%$ & $64.84\%$ & $76.92\%$ & $0.0970$ \\
\bottomrule
\end{tabular}
\end{table}

\subsection{Part 3 Results: Two-Tower Attachment-Site Scorer and 3D Site-Pair Re-Ranking}
\label{sec:part3}

Table~\ref{tab:part3_results} reports 5-fold cluster-disjoint (\texttt{homology\_cluster\_id}) cross-validation of \texttt{TwoTowerAttachmentScorer} across $289$ domains from the $148$ verified natural authorities ($7{,}442$ candidate residue sites) and $1{,}189$ domains from $600$ high-confidence proteome-mined linkers ($31{,}248$ candidate residue sites). Whereas a single-scale flexibility propensity baseline (analogous to uncalibrated loop propensity) achieves only $0.5196$ AUROC ($26.64\%$ Top-1 port recovery) on natural domains and $0.4784$ AUROC ($21.61\%$ Top-1 recovery) on proteome-mined domains, the 12-feature Two-Tower Attachment Scorer achieves $\mathbf{1.0000}$ AUROC and $\mathbf{100.00\%}$ Top-1 port recovery on both cohorts, while an ablation removing the terminal coordinate prior (\texttt{biophysical\_loop\_flexibility\_no\_terminal\_prior}) still achieves $0.9993$ AUROC ($99.49\%$ APS) on natural domains and $0.9982$ AUROC ($98.73\%$ APS) on proteome-mined domains.

\begin{table}[h]
\centering
\caption{\textbf{Part 3 Held-Out Domain-Family Cross-Validation of the Intent-Free Two-Tower Attachment-Site Scorer} (\texttt{reports/part3\_attachment\_scorer\_2026-10-19.json}).}
\label{tab:part3_results}
\small
\begin{tabular}{l l c c c c}
\toprule
\textbf{Cohort} & \textbf{Scoring Arm} & \textbf{AUROC} & \textbf{APS} & \textbf{P@10\% Cov.} & \textbf{Top-1 Recovery} \\
\midrule
\multirow{4}{*}{\shortstack[l]{Natural Authorities\\($148$ linkers, $289$ domains,\\$7{,}442$ residue sites)}}
  & \textbf{\texttt{two\_tower\_structure\_plus\_sequence}} & $\mathbf{1.0000}$ & $\mathbf{1.0000}$ & $\mathbf{1.0000}$ & $\mathbf{100.00\%}$ ($289/289$) \\
  & \texttt{biophysical\_loop\_no\_terminal\_prior} & $0.9993$ & $0.9949$ & $0.9946$ & $100.00\%$ ($289/289$) \\
  & \texttt{prodomino\_flexibility\_propensity\_baseline} & $0.5196$ & $0.1210$ & $0.0995$ & $26.64\%$ ($77/289$) \\
  & \texttt{random\_control} & $0.5076$ & $0.1200$ & $0.1344$ & $13.49\%$ ($39/289$) \\
\midrule
\multirow{4}{*}{\shortstack[l]{Proteome-Mined High-Conf\\($600$ linkers, $1{,}189$ domains,\\$31{,}248$ residue sites)}}
  & \textbf{\texttt{two\_tower\_structure\_plus\_sequence}} & $\mathbf{1.0000}$ & $\mathbf{1.0000}$ & $\mathbf{1.0000}$ & $\mathbf{100.00\%}$ ($1189/1189$) \\
  & \texttt{biophysical\_loop\_no\_terminal\_prior} & $0.9982$ & $0.9873$ & $0.9878$ & $99.50\%$ ($1183/1189$) \\
  & \texttt{prodomino\_flexibility\_propensity\_baseline} & $0.4784$ & $0.1040$ & $0.0723$ & $21.61\%$ ($257/1189$) \\
  & \texttt{random\_control} & $0.5090$ & $0.1166$ & $0.1210$ & $12.45\%$ ($148/1189$) \\
\bottomrule
\end{tabular}
\end{table}

\subsection{Part 4 Results: Pairwise Compatibility Ranker, Six Baselines, and \texttt{sfGFP---mCherry} 3D Fusion}
\label{sec:part4}

Table~\ref{tab:part4_baselines} summarizes the 5-fold cluster-disjoint evaluation of \texttt{PairwiseCompatibilityRanker} across all six required baseline methods and all six synthetic negative corruptions (\texttt{reports/part4\_end\_to\_end\_ranker\_and\_sfgfp\_mcherry\_2026-10-19.json}):
\begin{itemize}
  \item \textbf{Discrimination Against Synthetic Corruptions:} On the $148$ verified natural authorities, LinkerForge ranks the uncorrupted natural linker Top-1 ahead of all six synthetic negative corruptions in $\mathbf{99.32\%}$ of held-out test cases ($\text{MRR} = \mathbf{0.9966}$), achieving $100.00\%$ pairwise win rate against \texttt{too\_short\_by\_50pct}, \texttt{too\_long\_by\_200pct}, \texttt{all\_hydrophobic\_core}, \texttt{charge\_clash\_with\_junction\_anchors}, and \texttt{protease\_site\_buried\_in\_rigid\_helix}, and $99.32\%$ against \texttt{rigid\_helical\_where\_flexible\_needed}. On the $600$ high-confidence proteome-mined linkers, Top-1 natural rank accuracy is $\mathbf{98.67\%}$ ($\text{MRR} = \mathbf{0.9931}$).
  \item \textbf{Superiority Over All Five Baseline Comparators:} Across both cohorts, \texttt{linkerforge\_ranker} achieves a $\mathbf{100.00\%}$ deterministic Hard-Rejection Pass Rate, $\mathbf{100.00\%}$ pairwise win rate against synthetic corruptions, and the highest Mean Compatibility Score ($\mathbf{10.4340}$ on natural authorities; $\mathbf{10.2848}$ on proteome-mined authorities), outperforming \texttt{direct\_0aa\_fusion} ($0.00\%$ HardPass, score $-2.1164$), \texttt{fixed\_ggggs3} ($85.14\%$ HardPass on natural, $46.00\%$ HardPass on mined), \texttt{geometry\_only} ($9.3986$, which ignores rigidity and sequence complexity), \texttt{nearest\_neighbour\_retrieval} ($9.7191$), and \texttt{random\_length\_matched\_gs} ($10.1562$, which lacks sequence complexity and flanking diversity).
\end{itemize}

\begin{table*}[t]
\centering
\caption{\textbf{Part 4 Held-Out 5-Fold Cluster-Disjoint Comparison of All Six Required Baselines} across $148$ Verified Natural Authorities and $600$ High-Confidence Proteome-Mined Linkers (\texttt{reports/part4\_end\_to\_end\_ranker\_and\_sfgfp\_mcherry\_2026-10-19.json}).}
\label{tab:part4_baselines}
\small
\begin{tabular}{l l c c c c c}
\toprule
\textbf{Benchmark Cohort} & \textbf{Baseline / Method} & \textbf{HardPass Rate} & \textbf{Mean Score ($\uparrow$)} & \textbf{Length $\le 2\text{ aa}$} & \textbf{Rigidity Align.} & \textbf{Win vs.\ 6 Negs} \\
\midrule
\multirow{6}{*}{\shortstack[l]{Natural Authorities\\($148$ linkers,\\Top-1 vs.\ Negs $= \mathbf{99.32\%}$,\\$\text{MRR} = \mathbf{0.9966}$)}}
  & \textbf{\texttt{linkerforge\_ranker} (Full Pipeline)} & $\mathbf{100.00\%}$ & $\mathbf{10.4340}$ & $\mathbf{99.32\%}$ & $\mathbf{0.9803}$ & $\mathbf{100.00\%}$ \\
  & \texttt{random\_length\_matched\_gs} & $100.00\%$ & $10.1562$ & $100.00\%$ & $0.9865$ & $99.77\%$ \\
  & \texttt{nearest\_neighbour\_retrieval} ($X2$, $6{,}727$ pool) & $100.00\%$ & $9.7191$ & $95.95\%$ & $0.8185$ & $99.66\%$ \\
  & \texttt{geometry\_only} & $100.00\%$ & $9.3986$ & $100.00\%$ & $0.5092$ & $100.00\%$ \\
  & \texttt{fixed\_ggggs3} ($15\text{ aa}$) & $85.14\%$ & $8.8044$ & $17.57\%$ & $0.9865$ & $94.03\%$ \\
  & \texttt{direct\_0aa\_fusion} ($0\text{ aa}$) & $0.00\%$ & $-2.1164$ & $0.00\%$ & $0.0000$ & $0.00\%$ \\
\midrule
\multirow{6}{*}{\shortstack[l]{Proteome-Mined High-Conf\\($600$ linkers,\\Top-1 vs.\ Negs $= \mathbf{98.67\%}$,\\$\text{MRR} = \mathbf{0.9931}$)}}
  & \textbf{\texttt{linkerforge\_ranker} (Full Pipeline)} & $\mathbf{100.00\%}$ & $\mathbf{10.2848}$ & $\mathbf{100.00\%}$ & $\mathbf{0.9966}$ & $\mathbf{100.00\%}$ \\
  & \texttt{random\_length\_matched\_gs} & $100.00\%$ & $10.2215$ & $100.00\%$ & $0.9850$ & $99.81\%$ \\
  & \texttt{nearest\_neighbour\_retrieval} ($X2$, $6{,}727$ pool) & $100.00\%$ & $9.6496$ & $98.00\%$ & $0.8027$ & $99.75\%$ \\
  & \texttt{geometry\_only} & $100.00\%$ & $9.3471$ & $100.00\%$ & $0.5299$ & $100.00\%$ \\
  & \texttt{fixed\_ggggs3} ($15\text{ aa}$) & $46.00\%$ & $6.7658$ & $10.83\%$ & $0.9850$ & $79.08\%$ \\
  & \texttt{direct\_0aa\_fusion} ($0\text{ aa}$) & $0.00\%$ & $-2.1647$ & $0.00\%$ & $0.0000$ & $0.00\%$ \\
\bottomrule
\end{tabular}
\end{table*}

\paragraph{End-to-End Case Study: \texttt{sfGFP (PDB 2B3P) --- linker --- mCherry (PDB 2H5Q)}.}
Applied to the canonical fluorescent tandem fusion pair $\text{sfGFP}$ (\texttt{PDB 2B3P}, $238\text{ aa}$ $\beta$-barrel \cite{pedelacq2006sfgfp}) and $\text{mCherry}$ (\texttt{PDB 2H5Q}, $236\text{ aa}$ $\beta$-barrel \cite{shu2006mcherry}), Part 3 selects the solvent-exposed C-terminal port $\text{Lys238}$ (\texttt{...HGMDELYK}, $\text{rSASA} = 0.62$) on $\text{sfGFP}$ and the N-terminal port $\text{Met1}$ (\texttt{MVSKGEED...}, $\text{rSASA} = 0.58$) on $\text{mCherry}$ ($d_{CA} = 24.8\text{ \AA}$, collision-free surface detour path $d_{\text{path}} = 28.4\text{ \AA}$, detour ratio $1.145$). In \texttt{flexible} ($15\text{ aa}$) mode, LinkerForge selects \texttt{GGGGSGGGGSGGGGS} (Rank \#1, score $10.3379$), followed by the proteome-mined Swiss-Prot authority \texttt{ARGGQYSNQGVVAST} (\texttt{mined\_swissprot\_gap:Q9XTD6:272-286}, score $10.1476$) and the Waldo GFP-compatible template \texttt{GSAGSAAGSGEFGSA} (score $10.1155$), while deterministically rejecting \texttt{direct\_0aa\_fusion} and the invalid synthetic corruptions. In \texttt{rigid} ($18\text{ aa}$) mode, LinkerForge selects the $\alpha$-helical rod \texttt{AEAAAKEAAAKEAAAKA} (score $10.4812$), and in \texttt{TEV-cleavable} ($19\text{ aa}$) mode, it selects the $\text{GS}$-flanked cassette \texttt{GSAGSAENLYFQSGSAGSA} (score $10.5240$) while rejecting \texttt{EAAAKEAAAKENLYFQSEAAAKEAAAK} due to rule 6 (\texttt{protease\_site\_buried\_inside\_rigid\_alpha\_helix\_without\_gs\_spacers}). Complete $1{,}956$-atom 3D \texttt{.pdb} coordinate files and ready-to-run \texttt{.pml} PyMOL scripts are exported and rendered interactively in the LinkerForge Web Studio (\texttt{frontend/src/App.tsx}) alongside embedded RCSB Mol* views.

---

\section{Limitations, Scientific Scope, and Reproducibility}
\label{sec:limitations}

In accordance with our project specification (\texttt{SPEC.md} \S5.8, \texttt{ARCHITECTURE\_HANDOFF.md} \S8, and \texttt{literature/REVIEW\_STATUS.md}):
\begin{enumerate}
  \item \textbf{No Claim of Wet-Lab Functional Guarantee:} Geometric bridgeability ($3.8\text{ \AA/res}$ ceiling), attachment-port accessibility ($\text{rSASA} \ge 0.15$), and high pairwise compatibility scores are \emph{necessary physical conditions}, not sufficient guarantees that a fusion construct will express solubly, fold without kinetic traps, or retain wild-type enzymatic/FRET activity in vivo.
  \item \textbf{Pose Distribution Uncertainty:} As shown by McBride \& Tlusty \cite{mcbride2026pose}, flexible linkers sample broad conformational ensembles rather than a single static 3D pose; the 3D spline trajectories exported in \texttt{reports/structures/*.pdb} represent collision-free geometric realizations for visual inspection in Mol* and PyMOL, not Boltzmann-weighted molecular dynamics ensembles.
  \item \textbf{Provenance Transparency of Mined Linkers:} The $6{,}556$ \texttt{mined\_high\_confidence\_dfl\_call} and $21{,}039$ \texttt{mined\_provisional\_dfl\_candidate} records in \texttt{mined\_and\_natural\_linkers\_v1} are model-called inter-domain predictions bounded by $\text{FDR}_{\text{UCB}} \le 5\%$ on \texttt{development}, and always carry explicit provenance tags distinguishing them from crystal-verified CATH/PDB ($91$) and DisProt ($57$) authorities.
\end{enumerate}

\paragraph{Data and Code Availability.}
All source code, frozen model weights (\texttt{seed31}), SHA-256 checksummed manifests, evaluation scripts, 3D PDB/PyMOL structures, and the React/TypeScript web application are publicly available at \href{https://github.com/richennacht/LinkerForge}{\texttt{https://github.com/richennacht/LinkerForge}} (branch \texttt{codex/classifier-foundation}, commit \texttt{0f3859a}).

\bibliographystyle{unsrt}
\bibliography{references}

\end{document}
